\documentclass[11pt]{article}

\usepackage[utf8]{inputenc}
\usepackage[margin=1in]{geometry}
\usepackage{amsmath,amssymb,mathtools,enumitem,fancyhdr}
\usepackage[misc]{ifsym}
\usepackage{hyperref}

\setlength\textwidth{7in}
\setlength\parindent{0pt}
\oddsidemargin=-0.5cm
\textheight=665pt
\pagestyle{fancy}
\setlength{\headheight}{13.6pt}
\renewcommand{\headrulewidth}{0.1pt}
\renewcommand{\footrulewidth}{0.1pt}
\lhead{\scshape Pontificia Universidad Católica de Chile}
\rhead{\scshape IMC UC}
\cfoot{\thepage}

\newcommand{\solution}{\large\textbf{Solución}\normalsize}
\newcommand{\norm}[1]{\left\lVert #1\right\rVert}
\def\vec   #1{\mbox{\boldmath $#1$}{}}
\DeclareMathOperator{\dive}{div}
\def\R{\mathbb R}
\def\dx{\mbox{\(\,\mathrm{d}x\)}}
\def\ds{\mbox{\(\,\mathrm{d}s\)}}

\AtBeginDocument{\vspace*{-3.2\baselineskip}}

\begin{document}
    \begin{flushleft}
        \small
        \vspace{1pt}
        \textbf{IMT3130} \hfill
        \textbf{09/05/2026}
    \end{flushleft}
    \begin{center}
        \LARGE\textbf{Pauta Ayudantía 7}\\
        \vspace{3pt}
        \normalsize\textbf{Interpolación, elementos finitos compatibles e inf--sup}\\
        \normalsize Ayudante: Bastián Herrera \Letter{\hspace{1pt}: \texttt{biherrera@uc.cl}}
        \rule{\textwidth}{0.05mm}
    \end{center}

\section*{Problemas y soluciones}

\begin{enumerate}[label=\textbf{\arabic*.}, leftmargin=*]

    \item \textbf{Enunciado.} Para
    \[
        -u''+\eta u=f\quad\text{en }(0,1),\qquad u(0)=u(1)=0,
    \]
    con formulación débil $a_\eta(u,v)=F(v)$ en $V=H^1_0(0,1)$, compare los métodos de Galerkin conformes con $P_1$ y $P_2$. Pruebe ortogonalidad de Galerkin, derive Céa en norma de energía, estime los errores de energía por interpolación y use Aubin--Nitsche para obtener tasas en $L^2$.

    \textbf{Solución.}
    \begin{enumerate}[label=(\alph*)]
        \item La solución continua satisface
        \[
            a_\eta(u,v)=F(v)\qquad\forall v\in V,
        \]
        y la discreta
        \[
            a_\eta(u_h^{(r)},v_h)=F(v_h)\qquad\forall v_h\in V_h^{(r)}.
        \]
        Como $V_h^{(r)}\subset V$, restando ambas ecuaciones se obtiene
        \[
            a_\eta(u-u_h^{(r)},v_h)=0
            \qquad\forall v_h\in V_h^{(r)}.
        \]
        Sea $e_r=u-u_h^{(r)}$. Para cualquier $w_h\in V_h^{(r)}$,
        \begin{align*}
            \norm{e_r}_{a_\eta}^2
            &=a_\eta(e_r,e_r)\\
            &=a_\eta(e_r,u-u_h^{(r)})\\
            &=a_\eta(e_r,u-w_h)+a_\eta(e_r,w_h-u_h^{(r)})\\
            &=a_\eta(e_r,u-w_h)\\
            &\le \norm{e_r}_{a_\eta}\norm{u-w_h}_{a_\eta}.
        \end{align*}
        El término $a_\eta(e_r,w_h-u_h^{(r)})$ se anula porque $w_h-u_h^{(r)}\in V_h^{(r)}$ y vale la ortogonalidad de Galerkin.
        Si $e_r\neq0$, dividimos por $\norm{e_r}_{a_\eta}$ y tomamos ínfimo:
        \[
            \norm{u-u_h^{(r)}}_{a_\eta}
            \le
            \inf_{w_h\in V_h^{(r)}}\norm{u-w_h}_{a_\eta}.
        \]
        Esta es la estimación de Céa con constante $1$ en norma de energía.

        \item Tomando $w_h=\Pi_h^{(r)}u$,
        \[
        \begin{aligned}
            \norm{u-\Pi_h^{(r)}u}_{a_\eta}
            &=
            \left(
            |u-\Pi_h^{(r)}u|_{H^1(0,1)}^2
            +\eta\norm{u-\Pi_h^{(r)}u}_{L^2(0,1)}^2
            \right)^{1/2}\\
            &\le
            |u-\Pi_h^{(r)}u|_{H^1(0,1)}
            +\sqrt\eta\,\norm{u-\Pi_h^{(r)}u}_{L^2(0,1)}.
        \end{aligned}
        \]
        Para $P_1$,
        \[
        \begin{aligned}
            \norm{u-u_h^{(1)}}_{a_\eta}
            &\le
            C\left(h|u|_{H^2(0,1)}
            +\sqrt\eta\,h^2|u|_{H^2(0,1)}\right)\\
            &=
            C\left(h+\sqrt\eta\,h^2\right)|u|_{H^2(0,1)}
            =
            O(h)
            \quad(\eta\text{ fijo}).
        \end{aligned}
        \]
        Para $P_2$,
        \[
        \begin{aligned}
            \norm{u-u_h^{(2)}}_{a_\eta}
            &\le
            C\left(h^2|u|_{H^3(0,1)}
            +\sqrt\eta\,h^3|u|_{H^3(0,1)}\right)\\
            &=
            C\left(h^2+\sqrt\eta\,h^3\right)|u|_{H^3(0,1)}
            =
            O(h^2)
            \quad(\eta\text{ fijo}).
        \end{aligned}
        \]

        \item Sea $e_r=u-u_h^{(r)}$ y sea $z\in H^1_0(0,1)$ la solución dual
        \[
            a_\eta(w,z)=\int_0^1 w e_r\dx
            \qquad\forall w\in H^1_0(0,1).
        \]
        En forma fuerte, $z$ satisface
        \[
            -z''+\eta z=e_r\quad\text{en }(0,1),
            \qquad
            z(0)=z(1)=0.
        \]
        En 1D, para este problema elíptico, usamos la regularidad
        \[
            \norm{z}_{H^2(0,1)}\le C\norm{e_r}_{L^2(0,1)}.
        \]
        Entonces, para cualquier $z_h\in V_h^{(r)}$,
        \begin{align*}
            \norm{e_r}_{L^2(0,1)}^2
            &=
            a_\eta(e_r,z)\\
            &=
            a_\eta(e_r,z-z_h)
            \qquad\text{(ortogonalidad de Galerkin)}\\
            &\le
            \norm{e_r}_{a_\eta}\norm{z-z_h}_{a_\eta}.
        \end{align*}
        Como $V_h^{(r)}$ contiene funciones $P_1$, podemos escoger un interpolante lineal de $z$ y obtener
        \[
        \begin{aligned}
            \inf_{z_h\in V_h^{(r)}}\norm{z-z_h}_{a_\eta}
            &\le
            \norm{z-\Pi_h^{(1)}z}_{a_\eta}\\
            &\le
            C\left(h+\sqrt\eta\,h^2\right)\norm{z}_{H^2(0,1)}\\
            &\le
            Ch\norm{e_r}_{L^2(0,1)}
            \qquad(\eta\text{ fijo}).
        \end{aligned}
        \]
        Por tanto,
        \[
            \norm{e_r}_{L^2(0,1)}
            \le
            Ch\norm{e_r}_{a_\eta}.
        \]
        Combinando con las estimaciones de energía,
        \[
            \norm{u-u_h^{(1)}}_{L^2(0,1)}
            =
            O(h^2),
            \qquad
            \norm{u-u_h^{(2)}}_{L^2(0,1)}
            =
            O(h^3).
        \]
        La norma $L^2$ es más débil que la norma de energía, y el argumento dual permite ganar un orden adicional.
    \end{enumerate}

    \item \textbf{Enunciado.} Estudie dos triples de Ciarlet. Primero, en $\widehat K=[-1,1]$, use $\mathcal P=\mathbb P_2(\widehat K)$ y los grados de libertad $p(-1)$, $p(0)$, $p(1)$ para probar unisolvencia y construir la base local. Luego, en el triángulo de referencia, use
    \[
        RT_0(\widehat K)=\{\vec q(x,y)=(\alpha,\beta)+\gamma(x,y):\alpha,\beta,\gamma\in\R\}
    \]
    con grados de libertad dados por flujos normales en aristas. Pruebe unisolvencia y compare la continuidad global de $P_2$ y $RT_0$.

    \textbf{Solución.}
    \begin{enumerate}[label=(\alph*)]
        \item Sea $p\in\mathbb P_2([-1,1])$. Si
        \[
            p(-1)=p(0)=p(1)=0,
        \]
        entonces $p$ tiene tres raíces distintas. Como $\deg p\le2$, necesariamente $p\equiv0$. Además, $\dim\mathbb P_2=3$ y hay tres grados de libertad, luego el triple es unisolvente.

        La base local se construye imponiendo $\phi_i(\widehat x_j)=\delta_{ij}$. Con la fórmula de Lagrange,
        \[
            \phi_i(x)=\prod_{j\neq i}
            \frac{x-\widehat x_j}{\widehat x_i-\widehat x_j},
            \qquad
            \widehat x_1=-1,\quad \widehat x_2=0,\quad \widehat x_3=1.
        \]
        Así, las bases locales son
        \[
            \phi_1(x)=\frac{x(x-1)}{2},
            \qquad
            \phi_2(x)=1-x^2,
            \qquad
            \phi_3(x)=\frac{x(x+1)}{2}.
        \]
        Satisfacen $\phi_i(\widehat x_j)=\delta_{ij}$ para $\widehat x_1=-1$, $\widehat x_2=0$, $\widehat x_3=1$.

        \item Escribimos
        \[
            \vec q(x,y)=(\alpha,\beta)+\gamma(x,y).
        \]
        Tomamos las aristas
        \[
            e_1=\{x=0,\ 0\le y\le1\},\quad
            e_2=\{y=0,\ 0\le x\le1\},\quad
            e_3=\{x+y=1\},
        \]
        con normales exteriores
        \[
            \vec n_1=(-1,0),\qquad
            \vec n_2=(0,-1),\qquad
            \vec n_3=\frac{1}{\sqrt2}(1,1).
        \]
        Entonces
        \[
            \ell_1(\vec q)=\int_{e_1}\vec q\cdot\vec n_1\ds=-\alpha,
            \qquad
            \ell_2(\vec q)=\int_{e_2}\vec q\cdot\vec n_2\ds=-\beta.
        \]
        En la hipotenusa, parametrizando $(x,y)=(s,1-s)$, $0\le s\le1$, se tiene $\ds=\sqrt2\,ds$ y
        \[
            \vec q(s,1-s)\cdot\vec n_3\,\ds
            =
            \frac{\alpha+\beta+\gamma}{\sqrt2}\sqrt2\,ds.
        \]
        Luego
        \[
            \ell_3(\vec q)=\alpha+\beta+\gamma.
        \]
        Si los tres grados de libertad son cero, entonces $\alpha=0$, $\beta=0$ y $\gamma=0$. Por tanto $\vec q\equiv\vec0$. Como $\dim RT_0(\widehat K)=3$ y hay tres grados de libertad, el elemento es unisolvente.

        \item El elemento $P_2$ se ensambla igualando valores nodales compartidos. Por eso genera funciones globales continuas y conformes en $H^1$. En cambio, $RT_0$ se ensambla igualando los flujos normales sobre aristas compartidas. Por eso genera campos con continuidad normal y conformes en $H(\dive)$, pero no necesariamente funciones en $[H^1(\Omega)]^2$.
    \end{enumerate}

    \item \textbf{Enunciado.} Para el problema de transporte--reacción
    \[
        u'+u=f\quad\text{en }(0,1),\qquad u(0)=0,
    \]
    considere
    \[
        U=\{u\in H^1(0,1):u(0)=0\},\qquad
        V=L^2(0,1),\qquad
        a(u,v)=\int_0^1(u'+u)v\dx.
    \]
    Muestre que $a$ no es coerciva en norma $H^1$, pruebe una condición inf--sup continua y verifique que falla para $U_h=\operatorname{span}\{x,x^2\}$ y $V_h=\operatorname{span}\{1\}$.

    \textbf{Solución.}
    \begin{enumerate}[label=(\alph*)]
        \item Si $a$ fuera coerciva respecto de $\norm{\cdot}_{H^1(0,1)}$, existiría $\alpha>0$ tal que
        \[
            a(u,u)\ge\alpha\norm{u}_{H^1(0,1)}^2
            \qquad\forall u\in U.
        \]
        Pero
        \[
            a(u,u)=\int_0^1(u'+u)u\dx
            =
            \frac12u(1)^2+\norm{u}_{L^2(0,1)}^2,
        \]
        porque $u(0)=0$. Tomando $u_n(x)=\sin(n\pi x)$,
        \[
            a(u_n,u_n)=\frac12,
            \qquad
            \norm{u_n}_{H^1(0,1)}^2
            =
            \frac12+\frac{n^2\pi^2}{2}\to\infty.
        \]
        Por tanto, no puede existir una constante de coercividad en norma $H^1$.

        \item Sea $w=u'+u$. Si $u\neq0$, entonces $w\neq0$; de lo contrario, $u'+u=0$ y $u(0)=0$ implicarían $u\equiv0$. Elegimos $v=w\in L^2(0,1)$. Entonces
        \[
            \sup_{v\in L^2(0,1)\setminus\{0\}}
            \frac{a(u,v)}{\norm{v}_{L^2}}
            \ge
            \frac{\int_0^1w^2\dx}{\norm{w}_{L^2}}
            =
            \norm{u'+u}_{L^2}.
        \]
        Falta controlar $\norm{u}_{H^1}$ por $\norm{u'+u}_{L^2}$. La ecuación $u'+u=w$, $u(0)=0$, da
        \[
            u(x)=e^{-x}\int_0^x e^s w(s)\,ds.
        \]
        De esta fórmula y Cauchy--Schwarz,
        \[
        \begin{aligned}
            |u(x)|
            &\le
            \int_0^x e^{s-x}|w(s)|\,ds
            \le
            \int_0^1 |w(s)|\,ds
            \le
            \norm{w}_{L^2(0,1)}.
        \end{aligned}
        \]
        Por tanto,
        \[
            \norm{u}_{L^2(0,1)}
            \le
            C\norm{w}_{L^2(0,1)}.
        \]
        Además, como $u'=w-u$,
        \[
            \norm{u'}_{L^2(0,1)}
            \le
            \norm{w}_{L^2(0,1)}
            +\norm{u}_{L^2(0,1)}
            \le
            C\norm{w}_{L^2(0,1)}.
        \]
        Por tanto,
        \[
            \norm{u}_{H^1(0,1)}\le C\norm{u'+u}_{L^2(0,1)}.
        \]
        Así,
        \[
            \sup_{v\in V}
            \frac{a(u,v)}{\norm{v}_{L^2}}
            \ge
            C^{-1}\norm{u}_{H^1(0,1)}.
        \]
        Tomando ínfimo sobre $u\neq0$,
        \[
            \inf_{u\in U\setminus\{0\}}
            \sup_{v\in V\setminus\{0\}}
            \frac{a(u,v)}{\norm{u}_{H^1}\norm{v}_{L^2}}
            \ge C^{-1}>0.
        \]

        \item En el espacio discreto,
        \[
            U_h=\operatorname{span}\{x,x^2\},
            \qquad
            V_h=\operatorname{span}\{1\}.
        \]
        Para $u_h=\alpha x+\beta x^2$,
        \[
            a(u_h,1)
            =
            \int_0^1(u_h'+u_h)\dx
            =
            u_h(1)-u_h(0)+\int_0^1u_h\dx
            =
            \frac32\alpha+\frac43\beta.
        \]
        Eligiendo, por ejemplo,
        \[
            u_h(x)=8x-9x^2,
        \]
        se obtiene
        \[
            a(u_h,1)=0,
            \qquad
            u_h\neq0.
        \]
        Como todo $v_h\in V_h$ es múltiplo de $1$,
        \[
            \sup_{v_h\in V_h\setminus\{0\}}
            \frac{a(u_h,v_h)}{\norm{v_h}_{L^2}}
            =
            0.
        \]
        La constante inf--sup discreta es cero. El punto es que, para problemas no coercivos, la estabilidad continua no se hereda automáticamente al discretizar: hay que verificar una condición inf--sup discreta compatible con los espacios elegidos.
    \end{enumerate}

\end{enumerate}

\end{document}
